\section*{ЗАКЛЮЧЕНИЕ}
\addcontentsline{toc}{section}{ЗАКЛЮЧЕНИЕ}

Разработанная программная система оптимизирующего компилятора Common Lisp представляет собой полностью функциональный транслятор исходного кода на языке Common Lisp в байт-код, реализующий все основные возможности языка, а также виртуальную машину для выполнения исходных программ.

Основные результаты работы:

\begin{enumerate}
\item Изучен процесс компиляции языков программирования и функциональных языков в частности.
\item Изучены методы оптимизации компиляторов.
\item Спроектирован компилятор с оптимизациями.
\item Реализована компиляция S-выражения в промежуточную форму.
\item Реализована генерация кода ассемблера.
\item Реализована генерация байт-кода.
\item Реализована оптимизация промежуточной формы.
\end{enumerate}

Все требования, объявленные в техническом задании, были полностью реализованы, все задачи, поставленные в начале разработки проекта, были также решены.

Готовый рабочий проект представлен в виде программной системы компилятора Common Lisp. Исходный код находится в публичном доступе и внедрён в os-pi.
